% GENERATED FILE. Edit canonical lessons, then run npm run generate:books.
% canonical-source-hash: fnv1a64:030113aad26b0d5f
% canonical-lessons: MR-C367-pratha, MR-C367-rudhi, MR-C367-vicarsarni, MR-C367-pravrtti, MR-C367-samajik

\chapter{Customs and Ways of Thinking}
\label{ch:mr-customs-and-ways-of-thinking}
\hlchaptermodality{367}

\begin{chapteropening}
\textbf{By the end of this chapter:} \emph{I can talk about customs, ideologies and tendencies.}

Complete the last lesson of chapter 367: I can talk about customs, ideologies and tendencies.
\end{chapteropening}

\section[\texorpdfstring{\mr{प्रथा} (prathā)}{प्रथा (prathā)}]{\mr{प्रथा} (prathā) — a practice; a custom}

\label{lesson:MR-C367-pratha}



\begin{quote}
Before the new one: say the Marathi for renunciation, then the Marathi for welfare.
\end{quote}

\subsection*{You'll want to know: \mr{प्रथा}}
\textbf{\mr{प्रथा}} — \emph{prathā} — \textquotedblleft{}a practice; a custom\textquotedblright{}.

\begin{quote}
\mr{सणाला} \mr{नवे} \mr{कपडे} \mr{घालण्याची} \mr{प्रथा} \mr{आहे}.
\end{quote}

\begin{quote}
\emph{saṇālā nave kapḍe ghālaṇyācī prathā āhe.}
\end{quote}

\begin{quote}
There is a custom of wearing new clothes at festivals.
\end{quote}

Feminine.

X \mr{ची} \mr{प्रथा} \mr{आहे}, it is the custom to X.

\begin{cousinweb}
\textbf{\mr{प्रथा}} (\emph{prathā}) is a tatsama, taken from Sanskrit \emph{prathā}, \textquotedblleft{}spreading; fame\textquotedblright{} (Monier-Williams). Marathi uses it for a practice that has spread.
\end{cousinweb}

\begin{grammarlens}[title={What you've built}]
One more word for giving reasons and opinions.
\end{grammarlens}

\subsection*{Your turn}
\begin{itemize}
\raggedright
  \item \emph{Say it:} \emph{prathā}
  \item \emph{Say it:} \emph{prathā}, once more
  \item \emph{Say it:} \emph{saṇālā nave kapḍe ghālaṇyācī prathā āhe.} — There is a custom of wearing new clothes at festivals.
  \item \emph{Recall:} say the Marathi for renunciation, then the Marathi for welfare, then say \emph{prathā} again
\end{itemize}

\begin{tcolorbox}[breakable,colback=teal!4,colframe=teal!35!black,title={Before you move on}]
What does \mr{प्रथा} mean? (\textquotedblleft{}A practice; a custom\textquotedblright{}.) Say it once more.
\end{tcolorbox}
\section[\texorpdfstring{\mr{रूढी} (rūḍhī)}{रूढी (rūḍhī)}]{\mr{रूढी} (rūḍhī) — a convention; a tradition}

\label{lesson:MR-C367-rudhi}



\begin{quote}
Before the new one: say the Marathi for welfare, then the Marathi for a practice.
\end{quote}

\subsection*{You'll want to know: \mr{रूढी}}
\textbf{\mr{रूढी}} — \emph{rūḍhī} — \textquotedblleft{}a convention; a tradition\textquotedblright{}.

\begin{quote}
\mr{काही} \mr{जुन्या} \mr{रूढी} \mr{आता} \mr{बदलत} \mr{आहेत}.
\end{quote}

\begin{quote}
\emph{kāhī junyā rūḍhī ātā badlat āhet.}
\end{quote}

\begin{quote}
Some old conventions are now changing.
\end{quote}

Feminine.

Plural \mr{रूढी}; \mr{रूढ}, current, established.

\begin{cousinweb}
\textbf{\mr{रूढी}} (\emph{rūḍhī}) is a tatsama, taken from Sanskrit \emph{rūḍhi}, \textquotedblleft{}rise, increase\textquotedblright{} (CDIAL 10802); the same entry lists an inherited Marathi \emph{rūḍh}, \textquotedblleft{}prevalence\textquotedblright{}.
\end{cousinweb}

\begin{grammarlens}[title={What you've built}]
One more word for giving reasons and opinions.
\end{grammarlens}

\subsection*{Your turn}
\begin{itemize}
\raggedright
  \item \emph{Say it:} \emph{rūḍhī}
  \item \emph{Say it:} \emph{rūḍhī}, once more
  \item \emph{Say it:} \emph{kāhī junyā rūḍhī ātā badlat āhet.} — Some old conventions are now changing.
  \item \emph{Recall:} say the Marathi for welfare, then the Marathi for a practice, then say \emph{rūḍhī} again
\end{itemize}

\begin{tcolorbox}[breakable,colback=teal!4,colframe=teal!35!black,title={Before you move on}]
What does \mr{रूढी} mean? (\textquotedblleft{}A convention; a tradition\textquotedblright{}.) Say it once more.
\end{tcolorbox}
\section[\texorpdfstring{\mr{विचारसरणी} (vicārsarṇī)}{विचारसरणी (vicārsarṇī)}]{\mr{विचारसरणी} (vicārsarṇī) — an ideology; a way of thinking}

\label{lesson:MR-C367-vicarsarni}



\begin{quote}
Before the new one: say the Marathi for a practice, then the Marathi for a convention.
\end{quote}

\subsection*{You'll want to know: \mr{विचारसरणी}}
\textbf{\mr{विचारसरणी}} — \emph{vicārsarṇī} — \textquotedblleft{}an ideology; a way of thinking\textquotedblright{}.

\begin{quote}
\mr{दोघांची} \mr{विचारसरणी} \mr{अगदी} \mr{वेगळी} \mr{आहे}.
\end{quote}

\begin{quote}
\emph{doghāñcī vicārsarṇī agdī vegḷī āhe.}
\end{quote}

\begin{quote}
The two have very different ways of thinking.
\end{quote}

Feminine.

\mr{विचार} + \mr{सरणी}, thought + path.

\begin{cousinweb}
\textbf{\mr{विचारसरणी}} (\emph{vicārsarṇī}) is a learned compound built from \mr{विचार}, \textquotedblleft{}thought\textquotedblright{}, and Sanskrit \emph{saraṇi}, \textquotedblleft{}path, way\textquotedblright{} (Monier-Williams).
\end{cousinweb}

\begin{grammarlens}[title={What you've built}]
One more word for giving reasons and opinions.
\end{grammarlens}

\subsection*{Your turn}
\begin{itemize}
\raggedright
  \item \emph{Say it:} \emph{vicārsarṇī}
  \item \emph{Say it:} \emph{vicārsarṇī}, once more
  \item \emph{Say it:} \emph{doghāñcī vicārsarṇī agdī vegḷī āhe.} — The two have very different ways of thinking.
  \item \emph{Recall:} say the Marathi for a practice, then the Marathi for a convention, then say \emph{vicārsarṇī} again
\end{itemize}

\begin{tcolorbox}[breakable,colback=teal!4,colframe=teal!35!black,title={Before you move on}]
What does \mr{विचारसरणी} mean? (\textquotedblleft{}An ideology; a way of thinking\textquotedblright{}.) Say it once more.
\end{tcolorbox}
\section[\texorpdfstring{\mr{प्रवृत्ती} (pravṛttī)}{प्रवृत्ती (pravṛttī)}]{\mr{प्रवृत्ती} (pravṛttī) — a tendency}

\label{lesson:MR-C367-pravrtti}



\begin{quote}
Before the new one: say the Marathi for a convention, then the Marathi for an ideology.
\end{quote}

\subsection*{You'll want to know: \mr{प्रवृत्ती}}
\textbf{\mr{प्रवृत्ती}} — \emph{pravṛttī} — \textquotedblleft{}a tendency\textquotedblright{}.

\begin{quote}
\mr{प्रश्न} \mr{टाळण्याची} \mr{त्याची} \mr{प्रवृत्ती} \mr{आहे}.
\end{quote}

\begin{quote}
\emph{praśna ṭāḷaṇyācī tyācī pravṛttī āhe.}
\end{quote}

\begin{quote}
He has a tendency to avoid questions.
\end{quote}

Feminine.

X \mr{करण्याची} \mr{प्रवृत्ती}, a tendency to do X.

\begin{cousinweb}
\textbf{\mr{प्रवृत्ती}} (\emph{pravṛttī}) is a tatsama, taken from Sanskrit \emph{pravṛtti}, \textquotedblleft{}moving forward; inclination, tendency\textquotedblright{} (Monier-Williams).
\end{cousinweb}

\begin{grammarlens}[title={What you've built}]
One more word for giving reasons and opinions.
\end{grammarlens}

\subsection*{Your turn}
\begin{itemize}
\raggedright
  \item \emph{Say it:} \emph{pravṛttī}
  \item \emph{Say it:} \emph{pravṛttī}, once more
  \item \emph{Say it:} \emph{praśna ṭāḷaṇyācī tyācī pravṛttī āhe.} — He has a tendency to avoid questions.
  \item \emph{Recall:} say the Marathi for a convention, then the Marathi for an ideology, then say \emph{pravṛttī} again
\end{itemize}

\begin{tcolorbox}[breakable,colback=teal!4,colframe=teal!35!black,title={Before you move on}]
What does \mr{प्रवृत्ती} mean? (\textquotedblleft{}A tendency\textquotedblright{}.) Say it once more.
\end{tcolorbox}
\section[\texorpdfstring{\mr{सामाजिक} (sāmājik)}{सामाजिक (sāmājik)}]{\mr{सामाजिक} (sāmājik) — social}

\label{lesson:MR-C367-samajik}



\begin{quote}
Before the new one: say the Marathi for an ideology, then the Marathi for a tendency.
\end{quote}

\subsection*{You'll want to know: \mr{सामाजिक}}
\textbf{\mr{सामाजिक}} — \emph{sāmājik} — \textquotedblleft{}social\textquotedblright{}.

\begin{quote}
\mr{ती} \mr{अनेक} \mr{सामाजिक} \mr{कामांत} \mr{भाग} \mr{घेते}.
\end{quote}

\begin{quote}
\emph{tī anek sāmājik kāmānt bhāg ghete.}
\end{quote}

\begin{quote}
She takes part in many social causes.
\end{quote}

Invariable; built on \mr{समाज}, society.

\begin{cousinweb}
\textbf{\mr{सामाजिक}} (\emph{sāmājik}) is a tatsama, taken from Sanskrit \emph{sāmājika}, \textquotedblleft{}relating to an assembly\textquotedblright{} (Monier-Williams). Marathi builds it on \mr{समाज}, society.
\end{cousinweb}

\begin{grammarlens}[title={What you've built}]
One more word for giving reasons and opinions.
\end{grammarlens}

\subsection*{Your turn}
\begin{itemize}
\raggedright
  \item \emph{Say it:} \emph{sāmājik}
  \item \emph{Say it:} \emph{sāmājik}, once more
  \item \emph{Say it:} \emph{tī anek sāmājik kāmānt bhāg ghete.} — She takes part in many social causes.
  \item \emph{Recall:} say the Marathi for an ideology, then the Marathi for a tendency, then say \emph{sāmājik} again
\end{itemize}

\begin{tcolorbox}[breakable,colback=teal!4,colframe=teal!35!black,title={Before you move on}]
What does \mr{सामाजिक} mean? (\textquotedblleft{}Social\textquotedblright{}.) Say it once more.
\end{tcolorbox}
